\section*{Sensitivity to hip force direction (frontal-plane angle sweep)}
\label{sec:alpha-front-sweep}

\paragraph{Motivation.}
The hip contact force is a dominant driver of proximal femur loading and is therefore central to mechanostat-based remodeling simulations.
Instrumented hip implants provide in vivo contact force components, but their directions are time-dependent and vary across subjects and activities \cite{Bergmann2016_StandardizedLoadsHipImplants,OrthoLoadDatabase}.
Even for a fixed activity (walking), inter-individual peak-force directions vary on the order of $\sim 7^\circ$--$13^\circ$ in the frontal plane and $\sim 6^\circ$--$21^\circ$ in the sagittal plane \cite{Bergmann2016_StandardizedLoadsHipImplants}.
In addition, translating measured or musculoskeletal-modelled joint contact forces into a specific femur coordinate system introduces further uncertainty due to landmark identification and model parameter estimation \cite{Valente2014_MSKUncertainty}.
Because density adaptation integrates loading over long time horizons, it is important to quantify how such directional uncertainty propagates to spatial density features.

\paragraph{Sweep definition and objective.}
We probe directional sensitivity by applying a constant offset $\Delta\alpha_{\mathrm{front}}$ (in degrees) to the hip contact force frontal-plane angle $\alpha_{\mathrm{front}}$ in \emph{all} femur loading cases (heel strike, mid-stance, toe-off, stair climb), while leaving force magnitudes and muscle forces unchanged.
This isolates the effect of the hip force line-of-action in the frontal plane from other modelling choices.
We sweep $\Delta\alpha_{\mathrm{front}}\in[-15^\circ,\,25^\circ]$ (10 values), which spans the in vivo walking mean direction reported by Bergmann et al.\ ($\alpha_{\mathrm{front}}\approx 17^\circ$) as well as the ISO test direction mismatch (in vivo $\approx 15^\circ$--$24^\circ$ vs ISO $\approx 10^\circ$ in the frontal plane) \cite{Bergmann2016_StandardizedLoadsHipImplants}.
All runs use the accelerated femur benchmark parameter set (\textbf{StiffSetup}; Appendix Table~\ref{tab:param_setup_femur}) over $T=500$ days with identical numerical parameters.

\paragraph{Metrics for directional effects.}
Directional changes primarily redistribute bending/compression between medial and lateral cortices; we therefore use both (i) similarity-to-template metrics and (ii) medial--lateral redistribution metrics:
\begin{enumerate}
  \item \textbf{Similarity to CT template (population reference).}
  As in \S\ref{sec:femur-comparison}, we map the population-template CT density to the FE mesh (IDW interpolation) and rescale it to $[\rho_{\min},\rho_{\max}]$.
  We report nodal RMSE/MAE and Pearson correlation between $\rho_{\Delta\alpha}(\cdot,T)$ and $\rho_{\mathrm{CT}}$.
  These metrics are used diagnostically to compare sweep outcomes relative to the mapped population-template reference rather than as a claim of clinical accuracy.
  \item \textbf{Medial--lateral center-of-density.}
  Let $z(\mathbf{x})$ denote the medial axis coordinate in the femur coordinate system (FCS; Fig.~\ref{fig:femur_css}).
  We quantify redistribution by the mass-weighted first moment
  \[
  \bar z_{\rho}(\Delta\alpha)\;:=\;\frac{\int_\Omega \rho_{\Delta\alpha}(\mathbf{x},T)\,z(\mathbf{x})\,\mathrm{d}V}{\int_\Omega \rho_{\Delta\alpha}(\mathbf{x},T)\,\mathrm{d}V},
  \]
  and compare it against the corresponding CT value $\bar z_{\mathrm{CT}}$ computed on the same FE mesh.
  We additionally report the mass-weighted spread (standard deviation) $\sigma_z$ computed from the second moment to distinguish rigid shifts from broadening/narrowing of the density distribution.
  \item \textbf{Spatial sensitivity map.}
  To localize where density is most sensitive to the assumed hip direction, we compute a pointwise regression slope $\partial\rho/\partial\alpha_{\mathrm{front}}$ across the sweep and export it as a field for visualization.
\end{enumerate}

\paragraph{Hypothesis.}
We expect that increasing $\Delta\alpha_{\mathrm{front}}$ shifts remodeling stimulus toward the medial cortex (higher compressive load share) and thus increases medial density features, while decreasing $\Delta\alpha_{\mathrm{front}}$ shifts features laterally.
If the model is robust with respect to loading-direction uncertainty, similarity-to-template metrics should remain comparatively stable over a plausible range of $\Delta\alpha_{\mathrm{front}}$; conversely, strong variation would indicate that load direction is a sensitive modeling assumption that should be accounted for in robustness and uncertainty analyses.

\begin{figure}[htbp]
  \centering
  \safeincludegraphics[width=1\linewidth]{alpha_front_summary.png}
  \caption{Sensitivity to hip force frontal-plane angle offset $\Delta\alpha_{\mathrm{front}}$.
  (a) CT agreement (RMSE) versus angle offset.
  (b) Pattern agreement (Pearson correlation).
  (c) Medial--lateral redistribution (shift in center-of-density $\bar z_\rho - \bar z_{\mathrm{CT}}$).
  (d) Global density evolution over time.
  Shaded bands indicate the CT reference variability ($\pm 2\sigma$).
  }
  \label{fig:alpha_front_summary}
\end{figure}

\begin{figure}[htbp]
  \centering
  \safeincludegraphics[width=1\linewidth]{alpha_front_rho.png}
  \caption{Frontal-plane density distributions for varying hip force angles.
  The series illustrates the medial redistribution of bone density as the contact force angle shifts (left to right).
  Composite view from ParaView screenshots.}
  \label{fig:alpha_front_rho}
\end{figure}
